\chapter{Kalman Filter}
\label{cap:kalman}
% ── index ───────────────────────────────────────────────────────
\index{Kalman filter|textbf}
\index{kalman@\texttt{kalman()}|textbf}
\index{Rauch-Tung-Striebel smoother|textbf}
\index{RTS|see{Rauch-Tung-Striebel smoother}}
\index{state-space model|textbf}
\index{local level, model}
\index{local linear trend, model}
\index{unobservable state}
\index{prediction log-likelihood}

% ─────────────────────────────────────────────────────────────────────────────
\section{Hidden state and noisy observation}

The Kalman filter solves a different problem from all previous models:
estimating a \emph{latent} variable (the ``state'') that is not directly
observed, based on noisy measurements over time.

The general state-space model is:
\begin{align}
  y_t   &= H\, s_t + \varepsilon_t, \quad \varepsilon_t \sim N(0, R) \label{eq:obs}\\
  s_t   &= F\, s_{t-1} + R_s\, \eta_t, \quad \eta_t \sim N(0, Q) \label{eq:trans}
\end{align}
Equation \eqref{eq:obs} is the \textbf{observation equation}: $y_t$ is what is
measured, $s_t$ is the latent state and $\varepsilon_t$ is measurement noise. The
equation \eqref{eq:trans} is the \textbf{transition equation}: the state evolves
according to a linear dynamic perturbed by $\eta_t$.

\subsection*{Pre-defined models}

Hayashi offers two ready-made models for scalar series:

\medskip
\textbf{Local Level (LL):}
\[
  y_t = \mu_t + \varepsilon_t, \quad
  \mu_t = \mu_{t-1} + \eta_t
\]
The state $\mu_t$ is a random walk — the level drifts over time.
$\sigma_\varepsilon$ controls observation noise; $\sigma_\eta$ controls
the rate at which the level changes.

\medskip
\textbf{Local Linear Trend (LLT):}
\[
  y_t = \mu_t + \varepsilon_t, \quad
  \mu_t = \mu_{t-1} + \nu_{t-1} + \eta_t, \quad
  \nu_t = \nu_{t-1} + \zeta_t
\]
The state is two-dimensional: level $\mu_t$ and slope $\nu_t$.
The slope varies smoothly over time.

\subsection*{Kalman Filter (forward pass)}

The filter processes observations recursively:

\begin{enumerate}
  \item \textbf{Prediction}: $\hat{s}_{t|t-1} = F\,\hat{s}_{t-1|t-1}$,\quad
    $P_{t|t-1} = F P_{t-1|t-1} F' + R_s Q R_s'$

  \item \textbf{Innovation}: $v_t = y_t - H\,\hat{s}_{t|t-1}$,\quad
    $S_t = H P_{t|t-1} H' + R$

  \item \textbf{Kalman Gain}: $K_t = P_{t|t-1} H' S_t^{-1}$

  \item \textbf{Update}: $\hat{s}_{t|t} = \hat{s}_{t|t-1} + K_t v_t$,\quad
    $P_{t|t} = (I - K_t H) P_{t|t-1}$
\end{enumerate}

The result $\hat{s}_{t|t}$ is the \emph{filtered state}: optimal estimate of the
state using $y_1, \ldots, y_t$ (information up to $t$).

\subsection*{Rauch-Tung-Striebel Smoother (backward pass)}

After the filter, the RTS traverses the series backward and incorporates
future information:
\[
  L_t = P_{t|t}\, F'\, P_{t+1|t}^{-1}, \quad
  \hat{s}_{t|T} = \hat{s}_{t|t} + L_t\bigl(\hat{s}_{t+1|T} - \hat{s}_{t+1|t}\bigr)
\]
The result $\hat{s}_{t|T}$ is the \emph{smoothed state}: uses \emph{all}
the sample $y_1, \ldots, y_T$. The variance of the smoothed state is always less than or equal
to that of the filtered state — future information reduces uncertainty.

% ─────────────────────────────────────────────────────────────────────────────
\section{DGP Verification}

DGP: true level $\mu_t = 10 + 0{,}1\,t$ (linear trend) observed
with noise $\varepsilon_t \sim N(0, 1)$. $T = 150$.

\begin{lstlisting}[caption={Kalman Filter: LL and LLT}]
seed(42)
generate df mu = 10.0 + 0.1 * t
generate df y  = mu + noise       // noise ~ N(0,1)

// Local Level (sigma specified)
kalman(df, y, model="ll", sigma_obs=1.0, sigma_state=0.1)

// Local Linear Trend (automatic sigmas)
kalman(df, y, model="llt")
\end{lstlisting}

\subsection*{Output — Local Level}

\begin{lstlisting}[language={}, basicstyle=\ttfamily\small, frame=none,
  numbers=none, backgroundcolor=\color{codbg}]
Kalman (ll):  T=150  loglik=-278.1966  σ_obs=1.0000  σ_state=0.1000
  → y_filtered, y_smoothed added to df
\end{lstlisting}

\subsection*{Output — Local Linear Trend}

\begin{lstlisting}[language={}, basicstyle=\ttfamily\small, frame=none,
  numbers=none, backgroundcolor=\color{codbg}]
Kalman (llt):  T=150  loglik=-219.2354  σ_obs=0.9663  σ_state=0.0966  σ_slope=0.0097
  → y_filtered, y_smoothed, y_slope_filtered, y_slope_smoothed added to df
\end{lstlisting}

The LLT log-likelihood ($-219{,}2$) is substantially higher than the LL's
($-278{,}2$) because the DGP has a linear trend — the model with a free slope
better captures the data structure.

\subsection*{Comparison with the true state}

\begin{lstlisting}[language={}, basicstyle=\ttfamily\small, frame=none,
  numbers=none, backgroundcolor=\color{codbg}]
First 8 observations (LLT):

  t    mu     y       y_smoothed   y_slope_smoothed
  1   10.10  10.19      9.74           0.13
  2   10.20  10.35      9.87           0.13
  3   10.30  10.77      9.99           0.13
  4   10.40  10.07     10.10           0.13
  5   10.50   8.89     10.21           0.13
  6   10.60  10.31     10.33           0.12
  7   10.70  11.52     10.46           0.12
  8   10.80  12.01     10.56           0.12

summarize y_smoothed:  mean=17.49  sd=4.35
summarize mu:          mean=17.55  sd=4.34
\end{lstlisting}

The smoothed state recovers the true level with remarkable accuracy: mean 17{,}49
vs 17{,}55. The smoothed slope (\texttt{y\_slope\_smoothed} $\approx
0{,}13$) is close to the true value of $0{,}1$. The raw series $y$ oscillates
by $\pm 2$ around $\mu$; the smoother eliminates this noise.

\begin{nota}
  \textbf{Filtered vs smoothed}: \texttt{y\_filtered} uses only
  $y_1, \ldots, y_t$ — it is causal and useful for real-time forecasting.
  \texttt{y\_smoothed} uses the entire $y_1, \ldots, y_T$ — it is retrospective,
  but optimal for historical analysis. At interior points, the smoothed state
  is always more accurate; at the last point $t=T$, the two coincide.
\end{nota}

% ─────────────────────────────────────────────────────────────────────────────
\section{Syntax}

\begin{lstlisting}
// Local Level
kalman(df, y, model="ll")
kalman(df, y, model="ll", sigma_obs=1.0, sigma_state=0.1)

// Local Linear Trend
kalman(df, y, model="llt")
kalman(df, y, model="llt", sigma_obs=0.5, sigma_state=0.05, sigma_slope=0.005)
\end{lstlisting}

\begin{center}
\begin{tabular}{lll}
\toprule
\textbf{Parameter} & \textbf{Default} & \textbf{Description} \\
\midrule
\texttt{model}       & \texttt{"ll"}         & \texttt{"ll"} or \texttt{"llt"} \\
\texttt{sigma\_obs}  & \texttt{sd(diff(y))/\(\sqrt{2}\)} & Standard deviation of observation noise \\
\texttt{sigma\_state}& \texttt{sigma\_obs × 0.1}         & Standard deviation of level noise \\
\texttt{sigma\_slope}& \texttt{sigma\_state × 0.1}       & Standard deviation of slope noise (LLT) \\
\midrule
\multicolumn{3}{l}{\textit{Columns added to DataFrame:}} \\
\texttt{\{var\}\_filtered}        & — & Filtered level (forward pass) \\
\texttt{\{var\}\_smoothed}        & — & Smoothed level (RTS) \\
\texttt{\{var\}\_slope\_filtered} & — & Filtered slope (LLT only) \\
\texttt{\{var\}\_slope\_smoothed} & — & Smoothed slope (LLT only) \\
\bottomrule
\end{tabular}
\end{center}

\begin{center}
\begin{tabular}{lll}
\toprule
\textbf{Model} & \textbf{Use when} & \textbf{Relation} \\
\midrule
LL  & Level drift without clear trend    & Local level UCM (ch.~\ref{cap:filtros}) \\
LLT & Time-varying trend                 & Generalized stochastic HP filter \\
\bottomrule
\end{tabular}
\end{center}

\begin{nota}
  Hayashi uses the RTS smoother with fixed matrices specified by the
  user. For joint estimation of variance parameters
  ($\sigma_\varepsilon$, $\sigma_\eta$) via maximum likelihood,
  use the \texttt{ucm()} command (ch.~\ref{cap:filtros}), which runs
  the EM algorithm on the state-space representation.
\end{nota}
